\documentclass{article}
\usepackage[utf8]{inputenc}
\usepackage{a4wide}

\begin{document}

\section*{ประตูห้องลับ}

มีประตูบานหนึ่งสำหรับเข้าห้องเก็บสมบัติ

และมีคน\texttt{N}คนที่มีหมายเลขตั้งแต่\texttt{0}ถึง\texttt{N\ -\ 1}แต่ละคนสามารถเข้าหรือออกทางประตูได้หนึ่งครั้ง

โดยใช้เวลาหนึ่งวินาที\\



คุณจะได้รับ\texttt{arrival}เป็นอาร์เรย์จำนวนเต็มแบบไม่ลดลงขนาด\texttt{N}โดยที่\texttt{arrival{[}i{]}}คือเวลาที่มาถึงประตูของคนที่\texttt{i}และคุณยังได้รับอาร์เรย์\texttt{state}ขนาด\texttt{N}โดยที่\texttt{state{[}i{]}}จะเป็น\texttt{0}หากคนที่\texttt{i}ต้องการเข้าเดินประตู

หรือ 1 หากเขาต้องการเดินออกประตู หากมีคนตั้งแต่ 2

คนขึ้นไปต้องการใช้ประตูพร้อมกันให้ปฏิบัติตามกฎต่อไปนี้:\\



\begin{itemize}

\item

  หากไม่ได้ใช้ประตูในวินาทีก่อนหน้า ผู้ที่ต้องการออกจะได้ไปก่อน

\item

  หากมีการใช้ประตูในการเข้าในวินาทีก่อนหน้า ผู้ที่ต้องการเข้าไปก่อน

\item

  หากมีการใช้ประตูเพื่อออกในวินาทีก่อนหน้า ผู้ที่ต้องการออกจะไปก่อน

\item

  ถ้าต้องการไปในทิศทางเดียวกันหลายคน ให้คนที่มีดัชนีน้อยที่สุดไปก่อน

\end{itemize}



ให้ส่งคำตอบ\texttt{answer}เป็นอาร์เรย์ขนาด\texttt{N}โดยที่\texttt{answer{[}i{]}}คือวินาทีที่คนที่\texttt{i}ข้ามประตู



\hfill\break



\textbf{หมายเหตุ:}



\begin{enumerate}

\item

  ในแต่ละวินาที ให้มีได้เพียง 1 คนที่ข้ามประตู

\item

  แต่ละคนมาถึงประตูจะสามารถรอนิ่งๆ ได้โดยไม่ต้องเข้าหรือออกประตู

  เพื่อเป็นไปตามกติกาข้างบน

\end{enumerate}



\textbf{Example}



{\def\LTcaptype{none} % do not increment counter

\begin{longtable}[]{@{}

  >{\raggedright\arraybackslash}p{(\linewidth - 4\tabcolsep) * \real{0.3333}}

  >{\raggedright\arraybackslash}p{(\linewidth - 4\tabcolsep) * \real{0.3333}}

  >{\raggedright\arraybackslash}p{(\linewidth - 4\tabcolsep) * \real{0.3333}}@{}}

\toprule\noalign{}

\begin{minipage}[b]{\linewidth}\raggedright

Input

\end{minipage} & \begin{minipage}[b]{\linewidth}\raggedright

Output

\end{minipage} & \begin{minipage}[b]{\linewidth}\raggedright

Explanation\\

\strut

\end{minipage} \\

\midrule\noalign{}

\endhead

\bottomrule\noalign{}

\endlastfoot

\begin{minipage}[t]{\linewidth}\raggedright

5\\

0 1 1 2 4\\

0 1 0 0 1\\

\strut

\end{minipage} & \begin{minipage}[t]{\linewidth}\raggedright

0 3 1 2 4\textbf{\hfill\break

}\strut

\end{minipage} & \begin{minipage}[t]{\linewidth}\raggedright

\textbf{arrival = {[}0,1,1,2,4{]}, state = {[}0,1,0,0,1{]}\\

\strut \\

}ในแต่ละวินาที เกิดเหตุการณ์ต่อไปนี้:\\

- ที่ t = 0: คนที่ 0 เป็นเพียงคนเดียวที่ต้องการเข้าไป ดังนั้นเขาจึงเข้าไปทางประตู\\

- ที่ t = 1: คนที่ 1 ต้องการออก และคนที่ 2 ต้องการเข้า

เนื่องจากประตูถูกใช้ในวินาทีก่อนหน้าในการเข้า คนที่ 2 จึงเข้าไป (กฎข้อ 2)\\

- ที่ t = 2: คนที่ 1 ยังคงต้องการออก และคนที่ 3 ต้องการเข้า

เนื่องจากประตูถูกใช้ในวินาทีก่อนหน้าในการเข้า คนที่ 3 จึงได้เข้าไป (กฎข้อ 3)\\

- ที่ t = 3: คนที่ 1 เป็นเพียงคนเดียวที่ต้องการออก ดังนั้นเขาจึงออกทางประตู\\

- ที่ t = 4: คนที่ 4 เป็นเพียงคนเดียวที่ต้องการออก ดังนั้นเขาจึงออกทางประตูเท่านั้น\\

\strut \\

\strut

\end{minipage} \\

\begin{minipage}[t]{\linewidth}\raggedright

3\\

0 0 0\\

1 0 1\\

\strut

\end{minipage} & \begin{minipage}[t]{\linewidth}\raggedright

0 2 1\\

\strut

\end{minipage} & \begin{minipage}[t]{\linewidth}\raggedright

\textbf{arrival = {[}0,0,0{]}, state = {[}1,0,1{]}\\

}ในแต่ละวินาที เกิดเหตุการณ์ต่อไปนี้:\\

- ที่ t = 0: คนที่ 1 ต้องการเข้า และคนที่ 0 และ 2 ต้องการออก

เนื่องจากประตูไม่ได้ถูกใช้ในวินาทีก่อนหน้า\\

ดังนั้นคนที่ต้องการออกจะได้ออก (กฎข้อ 1) และเนื่องจาก (กฎข้อ\\

4) คนที่ 0 มีดัชนีต่ำกว่า 2 ดังนั้นคนที่ 0 จะได้ออกเป็นคนแรก (กฎข้อ 3)\\

- ที่ t = 1: คนที่ 1 ต้องการเข้า และคนที่ 2 ต้องการออก

เนื่องจากประตูถูกใช้เพื่อออกในวินาทีก่อนหน้า ดังนั้น 2 จึงได้ออก\\

- ที่ t = 2: เหลือคนที่ 1 คนเดียว ดังนั้นจึงเข้าประตูไปได้\\

\strut

\end{minipage} \\

\end{longtable}

}



Hints:



1. ใช้ 2 queue เพื่อคนที่ต้องการจะเข้า และคนที่ต้องการจะออก กับเวลา



2. simulate กระบวนการข้างต้นในโจทย์ โดยใช้กฎทั้ง 4 ข้อกับคนที่ต้องการจะเข้าหรือออก



\subsection*{Input}
บรรทัดที่ 1 เป็นเลขจำนวนเต็ม คือค่า\texttt{N}



บรรทัดที่ 2 เป็นชุดเลขจำนวนเต็ม\texttt{N}ค่าสำหรับ\texttt{arrival}เว้นด้วยช่องว่าง



บรรทัดที่ 3 เป็นชุดเลขจำนวนเต็ม\texttt{N}ค่าสำหรับ\texttt{state}เว้นด้วยช่องว่าง



\subsection*{Output}
ชุดเลขจำนวนเต็ม\texttt{N}ค่าสำหรับ\texttt{answer}เว้นด้วยช่องว่าง 1 ช่อง\\



\end{document}
